\section{A linear-distance theorem for Expand--Accumulate}
\label{sec:ea-distance}

This section turns the exact first moment from
Theorem~\ref{thm:expand-accumulate-first-moment} into a closed asymptotic
distance theorem for the Bernoulli ensemble of \cite{EA}.  The proof retains
the exact two-state accumulator process.  It exposes two independent
constraints: sparse messages determine the required expander density, while
linear-weight messages determine the achievable relative distance.

Let
\begin{equation}
 \mathsf{h}(x):=-x\ln x-(1-x)\ln(1-x)
 \label{eq:binary-entropy-nats}
\end{equation}
denote binary entropy in nats.  Fix $\delta\in(0,1/2)$ and define
\begin{equation}
 I_\delta:=1-2\sqrt{\delta(1-\delta)}.
 \label{eq:ea-sparse-rate}
\end{equation}

\subsection{Spectral extraction of the accumulator enumerator}

Fix a message of weight $r$.  By
Lemma~\ref{lem:bernoulli-expander-enumerator}, its expanded word has
independent Bernoulli coordinates with parameter $q_r$ from
\eqref{eq:bernoulli-expander-parity}.  If $z$ marks a nonzero accumulator
output, the exact transition matrix is
\begin{equation}
 T_q(z):=
 \begin{pmatrix}
  1-q & qz\\
  q   & (1-q)z
 \end{pmatrix}.
 \label{eq:ea-weight-matrix}
\end{equation}
Rows index the current accumulator state, and columns index the next state.
The factor $z$ marks the next state because that state is also the current
output symbol.

For $q\in[0,1/2]$ and $z\in(0,1]$, put
\begin{equation}
 \lambda(q,z)
 :=\frac{(1-q)(1+z)
 +\sqrt{(1-q)^2(1-z)^2+4q^2z}}{2}.
 \label{eq:ea-spectral-radius}
\end{equation}
For $q>0$, define the lower-tail exponent
\begin{equation}
 J_\delta(q)
 :=\sup_{0<z\leq 1}
 \left\{\delta\ln z-\ln\lambda(q,z)\right\}.
 \label{eq:ea-tail-rate}
\end{equation}

\begin{lemma}[Exact spectral tail bound]
\label{lem:ea-spectral-tail}
Let $H$ be the weight of the accumulator output when its input consists of
$n$ independent Bernoulli variables with parameter $q\in(0,1/2]$.  Then
\begin{equation}
 \Pr[H\leq\lfloor\delta n\rfloor]
 \leq \sqrt{2}\exp\bigl(-nJ_\delta(q)\bigr).
 \label{eq:ea-spectral-tail}
\end{equation}
Consequently, the Bernoulli Expand--Accumulate ensemble satisfies
\begin{equation}
 \Pr\left[\exists\xv\neq 0:
   \hw{\xv\B\A}\leq\lfloor\delta n\rfloor\right]
 \leq
 \sqrt{2}\sum_{r=1}^{k}\binom{k}{r}
 \exp\bigl(-nJ_\delta(q_r)\bigr).
 \label{eq:ea-spectral-first-moment}
\end{equation}
\end{lemma}

\begin{proof}
Let $\ev_0=(1,0)^{\mathsf T}$ and $\boldsymbol{1}=(1,1)^{\mathsf T}$.  The exact
probability generating function is
\begin{equation}
 \Ebb[z^H]=\ev_0^{\mathsf T}T_q(z)^n\boldsymbol{1}.
 \label{eq:ea-output-pgf}
\end{equation}
This identity is also the transfer-matrix form of the accumulator enumerator
in \eqref{eq:accumulator-enumerator-recalled}.

Set $D_z:=\operatorname{diag}(1,\sqrt z)$.  The similar matrix
\begin{equation}
 D_zT_q(z)D_z^{-1}
 =
 \begin{pmatrix}
  1-q & q\sqrt z\\
  q\sqrt z & (1-q)z
 \end{pmatrix}
 \label{eq:ea-symmetric-matrix}
\end{equation}
is symmetric and positive semidefinite.  Its largest eigenvalue is
$\lambda(q,z)$.  Therefore
\begin{equation}
 \Ebb[z^H]
 \leq \lambda(q,z)^n\lVert D_z\boldsymbol{1}\rVert_2
 =\sqrt{1+z}\,\lambda(q,z)^n.
 \label{eq:ea-pgf-spectral-bound}
\end{equation}
For $z\leq 1$, the event $H\leq\lfloor\delta n\rfloor$ implies
$z^H\geq z^{\lfloor\delta n\rfloor}$.  Apply Markov's inequality to
$z^H$, use \eqref{eq:ea-pgf-spectral-bound}, and optimize over $z$ to obtain
\eqref{eq:ea-spectral-tail}.  Sum this bound over the
$\binom{k}{r}$ messages of each weight $r$ to obtain
\eqref{eq:ea-spectral-first-moment}.
\end{proof}

The next lemma records the two endpoints of the exponent.  The limit at zero
controls sparse messages.  The value at one half controls the bulk.

\begin{lemma}[Endpoint exponents]
\label{lem:ea-endpoint-exponents}
For every $\delta\in(0,1/2)$, the function $J_\delta$ is continuous and
strictly positive on $(0,1/2]$.  Moreover,
\begin{align}
 \liminf_{q\downarrow 0}\frac{J_\delta(q)}{q}
 &\geq I_\delta,
 \label{eq:ea-sparse-exponent-limit}\\
 J_\delta(1/2)
 &=\ln 2-\mathsf{h}(\delta).
 \label{eq:ea-dense-exponent}
\end{align}
\end{lemma}

\begin{proof}
For the sparse limit, set $z=\exp(-\theta q)$ for a fixed $\theta\geq 0$.
Expansion of \eqref{eq:ea-spectral-radius} at $q=0$ gives
\begin{equation}
 \ln\lambda(q,e^{-\theta q})
 =q\left(-1-\frac{\theta}{2}
   +\sqrt{1+\frac{\theta^2}{4}}\right)+O(q^2).
 \label{eq:ea-sparse-spectral-expansion}
\end{equation}
Substitution in \eqref{eq:ea-tail-rate}, followed by the choice
\begin{equation}
 \theta=\frac{1-2\delta}{\sqrt{\delta(1-\delta)}},
 \label{eq:ea-sparse-theta}
\end{equation}
gives \eqref{eq:ea-sparse-exponent-limit}.

At $q=1/2$, one has $\lambda(1/2,z)=(1+z)/2$.  The optimizer in
\eqref{eq:ea-tail-rate} is $z=\delta/(1-\delta)$.  Direct substitution gives
\eqref{eq:ea-dense-exponent}.

It remains to justify the qualitative claims.  At $z=1$, the objective in
\eqref{eq:ea-tail-rate} is zero.  Its derivative with respect to $\ln z$ is
$\delta-1/2<0$.  Moving $z$ below one therefore makes the objective positive.
Thus $J_\delta(q)>0$ for every $q>0$.  On any compact interval bounded away
from $q=0$, the objective tends uniformly to minus infinity as $z$ tends to
zero.  Its supremum may consequently be restricted to a compact interval of
$z$ values.  Continuity then follows from continuity of
\eqref{eq:ea-spectral-radius}.
\end{proof}

\subsection{Asymptotic theorem}

\begin{theorem}[Binary Expand--Accumulate distance]
\label{thm:ea-linear-distance}
Let $k=k(n)$ satisfy $k/n\to R$ for a constant $R\in(0,1)$.  Sample
$\B\gets\dist^{\mathrm{Ber}}_{k,n,p_n}$ with
\begin{equation}
 p_n=C\frac{\ln n}{n}
 \label{eq:ea-density}
\end{equation}
for a constant $C>0$, and let $\A$ be the binary accumulator.  Fix
$\delta\in(0,1/2)$.  If
\begin{align}
 C\left(1-2\sqrt{\delta(1-\delta)}\right)&>1,
 \label{eq:ea-sparse-condition}\\
 \mathsf{h}(\delta)&<(1-R)\ln 2,
 \label{eq:ea-gv-condition}
\end{align}
then
\begin{equation}
 \Pr\left[
   \min_{\xv\in\Fbb_2^k\setminus\{0\}}
   \hw{\xv\B\A}\leq\lfloor\delta n\rfloor
 \right]=o(1).
 \label{eq:ea-linear-distance-conclusion}
\end{equation}
In particular, with probability $1-o(1)$, the sampled generator is injective
and its code has relative minimum distance greater than $\delta$.
\end{theorem}

\begin{proof}
We bound \eqref{eq:ea-spectral-first-moment} in three message-weight ranges.
Write $q_r=(1-(1-2p_n)^r)/2$.

First consider sparse messages.  By
\eqref{eq:ea-sparse-exponent-limit} and
\eqref{eq:ea-sparse-condition}, there are constants
$\epsilon>0$ and $a_0>0$ such that
\begin{equation}
 \gamma:=C(I_\delta-\epsilon)(1-a_0)>1
 \label{eq:ea-sparse-margin}
\end{equation}
and $J_\delta(q)\geq(I_\delta-\epsilon)q$ for $q\leq a_0$.
For $rp_n\leq a_0$, the inequalities
\begin{equation}
 q_r\leq rp_n,
 \qquad
 q_r\geq\frac{1-e^{-2rp_n}}{2}
       \geq rp_n(1-a_0)
 \label{eq:ea-qr-sparse-bounds}
\end{equation}
give
\begin{equation}
 nJ_\delta(q_r)\geq\gamma r\ln n.
 \label{eq:ea-sparse-tail-sum-exponent}
\end{equation}
Since $k=O(n)$, the contribution of all $r\leq a_0/p_n$ is at most
\begin{equation}
 \sqrt{2}\sum_{r\geq 1}k^r n^{-\gamma r}=o(1).
 \label{eq:ea-sparse-tail-sum}
\end{equation}

Next consider intermediate weights.  If $rp_n>a_0$, then
\begin{equation}
 q_r\geq q_*:=\frac{1-e^{-2a_0}}{2}>0.
 \label{eq:ea-intermediate-q}
\end{equation}
Lemma~\ref{lem:ea-endpoint-exponents} implies
\begin{equation}
 j_*:=\min_{q\in[q_*,1/2]}J_\delta(q)>0.
 \label{eq:ea-intermediate-rate}
\end{equation}
Choose a constant $\eta\in(0,R/2)$ small enough that
\begin{equation}
 R\mathsf{h}(\eta/R)<j_*.
 \label{eq:ea-intermediate-eta}
\end{equation}
The standard Hamming-ball estimate gives
\begin{equation}
 \sum_{r\leq\eta n}\binom{k}{r}
 \leq\exp\left(nR\mathsf{h}(\eta/R)+o(n)\right).
 \label{eq:ea-intermediate-message-count}
\end{equation}
Thus the contribution of $a_0/p_n<r\leq\eta n$ is exponentially small.

Finally, suppose that $r>\eta n$.  Uniformly over these weights,
\begin{equation}
 0\leq\frac12-q_r
 =\frac{(1-2p_n)^r}{2}
 \leq\frac12e^{-2p_nr}
 \leq\frac12n^{-2C\eta}=o(1).
 \label{eq:ea-dense-q}
\end{equation}
Continuity from Lemma~\ref{lem:ea-endpoint-exponents} therefore gives
\begin{equation}
 \inf_{r>\eta n}J_\delta(q_r)
 =\ln 2-\mathsf{h}(\delta)-o(1).
 \label{eq:ea-dense-rate}
\end{equation}
There are at most $2^k$ messages in this range.  Their total contribution is
at most
\begin{equation}
 \sqrt{2}\exp\left(
 n\bigl(R\ln2+\mathsf{h}(\delta)-\ln2+o(1)\bigr)
 \right)=o(1)
 \label{eq:ea-dense-tail-sum}
\end{equation}
by \eqref{eq:ea-gv-condition}.  All three ranges vanish, so
\eqref{eq:ea-spectral-first-moment} proves
\eqref{eq:ea-linear-distance-conclusion}.
\end{proof}

Condition~\eqref{eq:ea-gv-condition} is exactly the binary random-linear-code
distance condition.  The accumulator therefore loses no asymptotic distance
in the linear-weight regime.  The additional density condition
\eqref{eq:ea-sparse-condition} is caused by low-weight messages.  For example,
at relative distance $\delta=0.05$, it requires
\begin{equation}
 C>\frac{1}{1-2\sqrt{0.05\cdot0.95}}\approx1.773.
 \label{eq:ea-delta-five-density}
\end{equation}
This constant is an asymptotic threshold.  A finite failure target still
requires evaluation of the exact first moment in
Theorem~\ref{thm:expand-accumulate-first-moment}.

For comparison, write $\beta=1/2-\delta$.  The Markov-chain analysis in
\cite[Theorem~3.10]{EA} requires
\begin{equation}
 C>\frac{1}{\beta^2}
 \quad\text{and}\quad
 R<\min\left\{
 \frac{2}{\ln2}\frac{1-e^{-1}}{1+e^{-1}}\beta^2,
 \frac{2}{e^2}
 \right\}.
 \label{eq:ea-prior-conditions}
\end{equation}
At $\delta=0.05$, its density condition is $C>4.938$, whereas
\eqref{eq:ea-sparse-condition} requires only $C>1.773$.  At length
$n=5\cdot2^{20}$, these asymptotic thresholds correspond to expected row
weights greater than $76.41$ and $27.43$, respectively.  For rate $R=0.2$,
the prior rate condition permits $\delta<0.1127$, while
\eqref{eq:ea-gv-condition} permits $\delta<0.2430$.  These comparisons concern
asymptotic thresholds, not a certified finite failure probability.

\subsection{Certified finite evaluation}
\label{sec:ea-finite-certificate}

We now certify one parameter set from the concrete regime of \cite{EA}.  The
calculation must cover every nonzero message weight.  Evaluating a numerical
optimizer at selected weights does not provide that coverage.

For $q\in[0,1/2]$ and $z\in(0,1]$, define
\begin{equation}
 M_n(q,z):=\ev_0^{\mathsf T}T_q(z)^n\boldsymbol{1}.
 \label{eq:ea-exact-mgf}
\end{equation}
For an individually evaluated weight $r$ and a fixed marker $z_r$, the exact
transfer matrix gives
\begin{equation}
 \binom{k}{r}\Pr[H_r\leq L]
 \leq\binom{k}{r}z_r^{-L}M_n(q_r,z_r).
 \label{eq:ea-certified-exact-term}
\end{equation}
Here $H_r$ denotes the accumulator-output weight for any fixed message of
weight $r$.  The probability is over the Bernoulli expander.

The remaining weights are evaluated in contiguous blocks.  Recall the
symmetric matrix $S_q(z):=D_zT_q(z)D_z^{-1}$ from
\eqref{eq:ea-symmetric-matrix}.  Its derivative satisfies
\begin{equation}
 \frac{\partial S_q(z)}{\partial q}
 =-
 \begin{pmatrix}1\\-\sqrt z\end{pmatrix}
 \begin{pmatrix}1&-\sqrt z\end{pmatrix}
 \preceq0.
 \label{eq:ea-spectral-monotonicity}
\end{equation}
Thus $\lambda(q,z)$ is nonincreasing in $q$.  Since $q_r$ is nondecreasing in
$r$, any block $[a,b]\subseteq[1,k/2]$ satisfies
\begin{equation}
 \sum_{r=a}^{b}\binom{k}{r}\Pr[H_r\leq L]
 \leq
 (b-a+1)e^{k\mathsf{h}(b/k)}
 \sqrt{1+z}\,z^{-L}\lambda(q_a,z)^n.
 \label{eq:ea-certified-block-term}
\end{equation}
Equation~\eqref{eq:ea-certified-block-term} uses
$\binom{k}{r}\leq e^{k\mathsf{h}(r/k)}$ and monotonicity of binary entropy on
$[0,1/2]$.  For a block above $k/2$, the same bound holds after replacing the
message-count factor by $2^k$.

The certificate generator chooses the markers $z_r$ and $z$ with an ordinary
floating-point optimizer.  These choices do not affect soundness.  The
verifier treats every stored marker as a fixed decimal number.  It recomputes
\eqref{eq:ea-certified-exact-term} and
\eqref{eq:ea-certified-block-term} with outward-rounded Arb ball arithmetic.
The verifier also rejects a certificate unless its ranges partition
$[1,k]$.

\begin{theorem}[Certified concrete EA bounds]
\label{thm:ea-certified-concrete}
Let $(s,w,U)$ be any row of the following table.  Set
$k:=2^s$, $n:=5k$, $p:=w/n$, and $L:=n/20$.
\begin{center}
\begin{tabular}{c|r|c|c}
$s$ & expected row weight $w$ & failure bound $U$ & $-\log_2 U$\\
\hline
$20$ & $50$ & $6.891818499\cdot10^{-7}$ & $20.468611$\\
$25$ & $56$ & $7.475946766\cdot10^{-7}$ & $20.351240$\\
$30$ & $62$ & $8.107513660\cdot10^{-7}$ & $20.234237$
\end{tabular}
\end{center}
For $\B\gets\dist^{\mathrm{Ber}}_{k,n,p}$,
\begin{equation}
 \Pr_{\B}\left[
   \exists\xv\in\Fbb_2^k\setminus\{0\}:
   \hw{\xv\B\A}\leq L
 \right]
 \leq U<2^{-20}.
 \label{eq:ea-certified-result}
\end{equation}
\end{theorem}

\begin{proof}
Each checked certificate evaluates weights $1$ through $128$ with
\eqref{eq:ea-certified-exact-term}.  The three certificates cover the
remaining weights with $211$, $299$, and $388$ instances of
\eqref{eq:ea-certified-block-term}, respectively.  The largest total
contribution from these blocks is less than
$6.359\cdot10^{-976}$.  The outward upper endpoint of each complete Arb sum
is less than the corresponding displayed value $U$.
Theorem~\ref{thm:expand-accumulate-first-moment} proves
\eqref{eq:ea-certified-result}.
\end{proof}

The weight-one contribution dominates all three certificates.  Repeating the
calculations at the adjacent row weights $49$, $55$, and $61$ gives the bounds
$1.211510\cdot10^{-6}$, $1.314181\cdot10^{-6}$, and
$1.425203\cdot10^{-6}$, respectively.  These adjacent parameters do not meet
the 20-bit target under this certificate.

For comparison, the $C=3$ row of Figure~8 in \cite{EA} uses expected row
weights $46.42$, $56.81$, and $67.21$ at these lengths.  Its extrapolated
failure bounds at relative distance $0.05$ are $0.0157$, $0.00794$, and
$0.00402$.  The certified bounds above use a slightly heavier first row, but
they prove a uniform $2^{-20}$ failure target.  At the two larger lengths,
they also use lower expected row weights.

\subsection{Exact-row finite evaluation}
\label{sec:ea-exact-row-certificate}

We next replace the Bernoulli expander by
$\dist^{\mathrm{row}}_{k,n,d}$.  The positive first moment in
\eqref{eq:row-expand-accumulate-positive-first-moment} evaluates small
message weights directly.  Two bounds cover the remaining weights.

For $p=d/n$, define
\begin{equation}
 \vartheta_{n,d}:=\binom{n}{d}p^d(1-p)^{n-d}.
 \label{eq:exact-row-conditioning-probability}
\end{equation}

\begin{lemma}[Conditioned-Bernoulli bound]
\label{lem:exact-row-conditioning}
Fix a message of weight $r$, and let $E$ be an event determined by its
expanded word.  If $\B_{\mathrm{row}}\gets
\dist^{\mathrm{row}}_{k,n,d}$ and
$\B_{\mathrm{Ber}}\gets\dist^{\mathrm{Ber}}_{k,n,d/n}$, then
\begin{equation}
 \Pr[E(\xv\B_{\mathrm{row}})]
 \leq \vartheta_{n,d}^{-r}
 \Pr[E(\xv\B_{\mathrm{Ber}})].
 \label{eq:exact-row-conditioning-bound}
\end{equation}
\end{lemma}

\begin{proof}
Under the Bernoulli distribution, condition each of the $r$ selected rows to
have weight $d$.  The conditioned rows are independent and uniform on the
weight-$d$ slice, so their joint distribution is the exact-row distribution.
The conditioning event has probability $\vartheta_{n,d}^r$.  Dropping that
event from the numerator of the conditional probability proves
\eqref{eq:exact-row-conditioning-bound}.
\end{proof}

The dense bound uses the normalized Krawtchouk value
\begin{equation}
 \kappa_d(j):=\frac{K_d^{(n)}(j)}{\binom{n}{d}}.
 \label{eq:normalized-krawtchouk}
\end{equation}
For a character of weight $j$, this value is the Fourier eigenvalue of one
uniform weight-$d$ row.

\begin{lemma}[Dense exact-row bound]
\label{lem:exact-row-dense}
Fix a message of weight $r$ and put
\begin{equation}
 a_r:=\frac{2rd(n-d)}{n(n-1)}.
 \label{eq:exact-row-dense-exponent}
\end{equation}
For $L<n/2$, its exact-row Expand--Accumulate image satisfies
\begin{equation}
 \Pr[\hw{\xv\B\A}\leq L]
 \leq
 \left(
  \frac{e^{n\mathsf{h}(L/n)}}{2^n}\,
  2(1+e^{-a_r})^n
 \right)^{1/2}.
 \label{eq:exact-row-dense-bound}
\end{equation}
\end{lemma}

\begin{proof}
Let $X$ be the intersection size of a uniform $d$-subset and a fixed
$j$-subset of $[n]$.  Then $X$ is hypergeometric and
$\kappa_d(j)=\Ebb[(-1)^X]$.  A hypergeometric random variable has the same
distribution as a sum of independent, not necessarily identically
distributed, Bernoulli variables~\cite{HuiPark14}.  Writing their parameters
as $p_i$ gives
\begin{equation}
 |\kappa_d(j)|^2
 =\prod_i(1-2p_i)^2
 \leq\exp\left(-4\sum_i p_i(1-p_i)\right)
 =\exp\left(-4\operatorname{Var}(X)\right).
 \label{eq:krawtchouk-variance-bound}
\end{equation}
The hypergeometric variance is
\begin{equation}
 \operatorname{Var}(X)
 =d\frac{j}{n}\left(1-\frac{j}{n}\right)\frac{n-d}{n-1}.
 \label{eq:hypergeometric-variance}
\end{equation}

Let $Y$ be the sum of the $r$ selected exact-weight rows.  Parseval's identity
for the binary cube gives the squared $L_2$ norm of the density of $Y$ relative
to the uniform distribution as
\begin{equation}
 \sum_{j=0}^{n}\binom{n}{j}|\kappa_d(j)|^{2r}.
 \label{eq:exact-row-density-l2}
\end{equation}
The summand is symmetric under $j\mapsto n-j$.  For $j\leq n/2$,
\eqref{eq:krawtchouk-variance-bound} and
\eqref{eq:hypergeometric-variance} give
\begin{equation}
 |\kappa_d(j)|^{2r}\leq e^{-a_rj}.
 \label{eq:exact-row-eigenvalue-profile}
\end{equation}
Thus \eqref{eq:exact-row-density-l2} is at most
$2(1+e^{-a_r})^n$.

The accumulator is invertible.  The preimage of its output ball therefore has
the same cardinality, at most $e^{n\mathsf{h}(L/n)}$.  Apply Cauchy--Schwarz
between its indicator and the density of $Y$ to obtain
\eqref{eq:exact-row-dense-bound}.
\end{proof}

\begin{theorem}[Certified exact-row EA bounds]
\label{thm:ea-certified-exact-row}
Let $(s,d,U)$ be any row of the following table.  Set $k:=2^s$, $n:=5k$, and
$L:=n/20$.  If $\B\gets\dist^{\mathrm{row}}_{k,n,d}$, then
\begin{equation}
 \Pr_{\B}\left[
  \exists\xv\in\Fbb_2^k\setminus\{0\}:
  \hw{\xv\B\A}\leq L
 \right]\leq U<2^{-20}.
 \label{eq:ea-certified-exact-row-result}
\end{equation}
\begin{center}
\begin{tabular}{c|r|c|c}
$s$ & exact row weight $d$ & failure bound $U$ & $-\log_2 U$\\
\hline
$20$ & $31$ & $3.321605945\cdot10^{-7}$ & $21.521615$\\
$25$ & $35$ & $2.548539646\cdot10^{-7}$ & $21.903825$\\
$30$ & $39$ & $2.796467641\cdot10^{-7}$ & $21.769891$
\end{tabular}
\end{center}
\end{theorem}

\begin{proof}
For $s=20,25,30$, respectively, the verifier evaluates message weights through
$7,16,48$ with the positive shell recurrence
\eqref{eq:row-shell-recurrence} and the positive accumulator tail
\eqref{eq:accumulator-slice-tail}.  It covers the interval up to $k/8-1$
with $225,304,373$ contiguous conditioned-Bernoulli blocks.  Each block uses
\eqref{eq:ea-certified-block-term} and pays the conditioning factor from
\eqref{eq:exact-row-conditioning-bound}.  Lemma~\ref{lem:exact-row-dense}
covers every weight from $k/8$ through $k$ at once.

The implementation evaluates the shell and conditioning probabilities with
outward-rounded Arb balls.  A floating-point optimizer selects the Chernoff
marker for each conditioning block, but the verifier treats the resulting
decimal as fixed and checks the bound with Arb.  The three dense contributions
are below $2^{-28{,}342}$, $2^{-2{,}082{,}004}$, and
$2^{-98{,}337{,}579}$, respectively.  Summing all three regimes gives the
displayed values.  The first-moment bound in
Theorem~\ref{thm:expand-accumulate-first-moment} proves
\eqref{eq:ea-certified-exact-row-result}.
\end{proof}

The exact-row table uses $19$, $21$, and $23$ fewer nonzero entries per row
than the Bernoulli table in Theorem~\ref{thm:ea-certified-concrete}.  In all
three exact-row certificates, the weight-one message term is the largest
directly evaluated term.

\subsection{Region-stratified left-regular evaluation}
\label{sec:ea-regular-certificate}

The global exact-row ensemble controls row weight but not the location of its
nonzero entries.  We now apply the region-stratified ensemble from
Section~\ref{sec:region-left-regular-expander}.  The accumulator retains one
state bit between regions, so a scalar regional weight enumerator is
insufficient.  Define the bivariate transfer matrix
\begin{equation}
 \mathbf U(X,Y):=
 \begin{pmatrix}
  1&XY\\
  X&Y
 \end{pmatrix}.
 \label{eq:ea-regular-bivariate-matrix}
\end{equation}
Rows index the current accumulator state, and columns index the next state.
The exponent of $X$ records input weight.  The exponent of $Y$ records output
weight.  For $u\in[0,\ell]$, define
\begin{align}
 \mathbf S^{\mathrm{acc}}_{\ell,u}(Y)
 &:=\frac{1}{\binom{\ell}{u}}[X^u]\mathbf U(X,Y)^\ell,
 \label{eq:ea-regular-slice-matrix}\\
 \mathbf R^{\mathrm{acc}}_{\ell,r}(Y)
 &:=\sum_{u=0}^{\ell}\rho^{(\ell)}_{r,u}
      \mathbf S^{\mathrm{acc}}_{\ell,u}(Y).
 \label{eq:ea-regular-region-matrix}
\end{align}

\begin{theorem}[Exact left-regular EA first moment]
\label{thm:ea-regular-first-moment}
Assume that $n=d\ell$.  Sample
$\B\gets\dist^{\mathrm{reg}}_{k,n,d}$ and set $\G:=\B\A$.  Let $\ev_0$ be
the unit column vector for accumulator state zero.  The expected number of
nonzero messages whose encoding has weight at most $L$ is
\begin{equation}
 \boxed{
 \mu_{\mathrm{EA}}^{\mathrm{reg}}(L)
 =\sum_{r=1}^{k}\binom{k}{r}\sum_{h=0}^{L}[Y^h]\,
   \ev_0^\mathsf{T}
   \mathbf R^{\mathrm{acc}}_{\ell,r}(Y)^d\boldsymbol{1}.
 }
 \label{eq:ea-regular-first-moment}
\end{equation}
This quantity upper-bounds the probability that the generated code has
minimum distance at most $L$.
\end{theorem}

\begin{proof}
Lemma~\ref{lem:region-left-regular-shell-law} makes the input in each region a
mixture of uniform Hamming slices.  Matrix
\eqref{eq:ea-regular-slice-matrix} averages the accumulator transfer over one
such slice.  Matrix multiplication retains the accumulator state at each
regional boundary.  Sum over regional shells, output weights, and nonzero
messages.  Markov's inequality gives the failure bound.
\end{proof}

The exact matrix coefficient is used only for small message weights.  Lemma
\ref{lem:regular-poisson-comparison} and the Bernoulli accumulator matrix
$T_q(z)$ give
\begin{equation}
 \Ebb_{\B}\!\left[z^{\hw{\xv\B\A}}\right]
 \leq\gamma_r^d\,
 \ev_0^\mathsf{T}T_{q^{\mathrm{reg}}_r}(z)^n\boldsymbol{1}.
 \label{eq:ea-regular-poisson-mgf}
\end{equation}
For a contiguous block $s\leq r\leq t\leq k/2$, use $gamma_t$ and
$q^{\mathrm{reg}}_s$.  The former is an upper bound because $\gamma_r$ is
nondecreasing.  The latter is an upper bound on the lower-tail moment because
the spectral radius in \eqref{eq:ea-spectral-radius} is nonincreasing in $q$.

For dense messages, let $V_n(L):=\sum_{h=0}^{L}\binom{n}{h}$.  The accumulator
is invertible.  Lemma~\ref{lem:regular-point-mass} therefore gives
\begin{equation}
 \Pr_{\B}[\hw{\xv\B\A}\leq L]
 \leq V_n(L)2^{d-n}(1+e^{-2r/\ell})^n.
 \label{eq:ea-regular-dense-bound}
\end{equation}

\begin{theorem}[Certified left-regular EA bounds]
\label{thm:ea-certified-regular}
For each row below, let $n:=d\ell$, sample
$\B\gets\dist^{\mathrm{reg}}_{k,n,d}$, and use the binary accumulator $\A$.
The displayed value $U$ satisfies
\begin{equation}
 \Pr_{\B}[\exists\xv\ne0:\ \hw{\xv\B\A}\leq L]
 \leq U<2^{-20}.
 \label{eq:ea-certified-regular-result}
\end{equation}
\begin{center}
\begin{tabular}{c|r|r|r|r|c|c}
rate &$k$ &$d$ &$\ell$ &$L$ &$U$ &$-\log_2U$\\
\hline
$1/5$ &$1{,}048{,}544$ &$31$ &$169{,}120$ &$262{,}136$
 &$3.108423898\cdot10^{-7}$ &$21.617313$\\
$1/2$ &$1{,}048{,}576$ &$64$ &$32{,}768$ &$230{,}718$
 &$8.982345581\cdot10^{-7}$ &$20.086404$
\end{tabular}
\end{center}
The rate-$1/2$ distance is
\[
 \frac{L}{n}=0.1100149155\ldots
 =0.9998823119\ldots\,\delta_{\mathrm{GV}}(1/2).
\]
\end{theorem}

\begin{proof}
The verifier \path{scripts/regular_ea_certificate.py} uses
\eqref{eq:ea-regular-first-moment} through message weights $17$ and $14$ in
the two rows.  The corresponding exact contributions are at most
$3.108394\cdot10^{-7}$ and $8.933170\cdot10^{-7}$.  In both cases, message
weight two gives the largest exact term.

The verifier next applies \eqref{eq:ea-regular-poisson-mgf} in $18$ and $13$
contiguous blocks.  Their total contributions are at most
$3.001904\cdot10^{-12}$ and $1.153645\cdot10^{-12}$.  Finally, it applies
\eqref{eq:ea-regular-dense-bound} in $3072$ and $3073$ positive blocks of width
at most $256$.  The dense contributions are below $2^{-700{,}000}$ and
$4.916442\cdot10^{-9}$, respectively.  Every bound is evaluated with
outward-rounded Arb balls at $192$-bit precision.  Summing the three regimes
proves \eqref{eq:ea-certified-regular-result}.
\end{proof}

At rate $1/5$, the regular degree $31$ matches the smallest certified global
exact-row weight in Theorem~\ref{thm:ea-certified-exact-row}; it does not
improve expander sparsity.  At rate $1/2$, the Bernoulli certificate at the
same $k,n,L$ requires expected row weight $75$ to reach the 20-bit target.  It
gives failure bound $8.174250904\cdot10^{-7}$.  The regular row weight $64$
reduces the number of expander edges by $14.67\%$.  This gain is smaller than
the wrapped-EC gain because every even-weight message returns the accumulator
to state zero at each regional boundary.  Message weight two is consequently
the limiting case.
